% titel: Blatt 2, A1: Eine neue Golfvariante (schiefer Wurf)
% skript: Kap. Kinematik, Projektilbewegung

% Wurfparabel, Schnittpunkt mit einer Wand, Reichweite.

Beim schiefen Wurf ohne Luftreibung überlagern sich eine gleichförmige
Bewegung in horizontaler und eine gleichmässig beschleunigte in
vertikaler Richtung: die Bahn ist eine Parabel.

\begin{werkzeugkasten}[Werkzeuge]
Mit $\vec g$ nach unten und Start im Ursprung:
\begin{equation}
  x(t)=v_0\cos\alpha\,t,\qquad
  y(t)=v_0\sin\alpha\,t-\tfrac12 g t^2.
\end{equation}
Ohne $t$ ergibt sich die \textbf{Wurfparabel}
\begin{equation}
  y(x)=x\tan\alpha-\frac{g\,x^2}{2\cos^2\alpha\,v_0^2}\;.
\end{equation}
\textbf{Reichweite} (Treffer $y=0$ beim Abflug, $y=0$ beim Aufschlag):
\begin{equation}
  R=\frac{v_0^2\,\sin(2\alpha)}{g}\;.
\end{equation}

\textbf{Energieerhaltung} (keine Reibung): Geschwindigkeit in gleicher
Hoehe wieder $v_0$; gilt allgemein für beliebige Haltepunkte der Bahn.
\end{werkzeugkasten}

\begin{werkzeugkasten}[Vorgehen]
\begin{itemize}
  \item Die Gr\"o\"sen $v_0$ und $\alpha$ stehen nur in der Parabelformel;
        einen gegebenen Punkt $(a,h)$ einsetzen.
  \item Den Abstand hinter der Wand als Reststück $R-a$ bestimmen.
  \item Für die Endgeschwindigkeit Energieerhaltung (oder Parabelsymmetrie)
        nutzen.
\end{itemize}
\end{werkzeugkasten}

\begin{achtung}[Typische Fallen]
\begin{itemize}
  \item Die Wand loest \emph{nicht} die Bewegung: der Ball durchquert das
        Loch und fliegt als dieselbe Parabel weiter.
  \item Reichweite $R$ heisst der Abstand zwischen Abwurf und R\"uckkehr auf
        $y=0$; der Auftreffpunkt hinter der Wand liegt bei $x=R$.
  \item $\sin(2\alpha)$ statt $\cos(2\alpha)$ verwenden.
\end{itemize}
\end{achtung}
